\chapter{Node configuration reference}%
\label{ch:bijlage-config}

This appendix lists the full configuration of the three deployed nodes as a technical
reference, complementing the summary given in Table~\ref{tab:poc-config}. Values were
read out using the Meshtastic command-line interface (\texttt{--info}) after the final
configuration of each node.
Quick note: WiFi doesn't need to be enabled, but with the node on the roof, this was easier to access the rangetest csv file.

\section{Radio and device configuration}%
\label{sec:bijlage-config-device}

\begin{table}[h]
  \centering
  \small
  \begin{tabular}{lccc}
    \toprule
    \textbf{Parameter} & \textbf{Ronse-Base} & \textbf{Relay} & \textbf{Mobile} \\
    \midrule
    Role                    & CLIENT\_BASE & ROUTER      & CLIENT \\
    Firmware                & 2.7.26.54e0d8d & 2.7.26.54e0d8d & 2.7.26.54e0d8d \\
    Hardware model          & Heltec V4    & Heltec V4   & Heltec V4 \\
    Region                  & EU\_868      & EU\_868     & EU\_868 \\
    Modem preset            & LONG\_FAST   & LONG\_FAST  & LONG\_FAST \\
    Bandwidth               & 250\,kHz     & 250\,kHz    & 250\,kHz \\
    Spreading factor        & 11           & 11          & 11 \\
    Coding rate             & 4/5          & 4/5         & 4/5 \\
    Hop limit               & 3            & 3           & 3 \\
    TX power                & 27\,dBm      & 27\,dBm     & 27\,dBm \\
    Boosted RX gain         & Enabled      & Enabled     & Enabled \\
    FEM/LNA mode            & Disabled     & Disabled    & Disabled \\
    Rebroadcast mode        & ALL          & ALL         & ALL \\
    \midrule
    Fixed position          & Yes          & Yes         & No (GPS via phone) \\
    Latitude                & 50.72316     & 50.65457    & --- \\
    Longitude               & 3.63217      & 3.39975     & --- \\
    Altitude                & 90\,m        & 140\,m      & --- \\
    \midrule
    Power saving            & Off          & On          & Off \\
    Bluetooth                & Enabled, no PIN & Enabled, no PIN & Enabled, no PIN \\
    WiFi                     & Enabled (bridge) & Disabled & Disabled \\
    Display screen timeout   & 60\,s        & 30\,s       & 60\,s \\
    \bottomrule
  \end{tabular}
  \caption[Radio and device configuration]{\label{tab:bijlage-config}Radio and device configuration of the three nodes.}
\end{table}

\section{Module configuration}%
\label{sec:bijlage-config-modules}

On top of the default settings, Meshtastic exposes a number of independent modules
that were configured per node according to its role. Table~\ref{tab:bijlage-modules}
lists the modules relevant to this deployment; modules not used in this thesis
(such as Store\&Forward, Detection Sensor, Paxcounter, Audio and Canned Message)
were left at their firmware defaults (disabled) on all nodes.

\begin{table}[h]
  \centering
  \small
  \begin{tabular}{lccc}
    \toprule
    \textbf{Module} & \textbf{Ronse-Base} & \textbf{Relay} & \textbf{Mobile} \\
    \midrule
    \multicolumn{4}{l}{\textit{Range Test}} \\
    Enabled                 & Yes (receiver) & Yes (receiver) & Yes (sender) \\
    Sender interval         & ---          & ---         & 60\,s \\
    CSV logging             & Enabled      & Disabled    & --- \\
    \midrule
    \multicolumn{4}{l}{\textit{Telemetry}} \\
    Device telemetry        & Enabled      & Enabled     & Enabled \\
    Update interval         & 1,800\,s      & 1,800\,s     & 1,800\,s \\
    Environment telemetry   & Disabled     & Disabled    & Disabled \\
    \midrule
    \multicolumn{4}{l}{\textit{Neighbor Info}} \\
    Enabled                 & Yes          & Yes         & Yes \\
    Update interval         & 21,600\,s      & 21,600\,s     & 21,600\,s \\
    \midrule
    \multicolumn{4}{l}{\textit{MQTT}} \\
    Enabled                 & No           & No          & No \\
    \bottomrule
  \end{tabular}
  \caption[Module configuration]{\label{tab:bijlage-modules}Module configuration of the three deployed nodes.}
\end{table}

\section{Channel and security configuration}%
\label{sec:bijlage-config-security}

All three nodes were migrated from the Meshtastic default public channel to a private
channel secured with a randomly generated 256-bit \gls{psk}, as discussed in
Section~\ref{sec:poc-security}. The channel name and \gls{psk} are identical across all three
nodes so that they can communicate; the \gls{psk} itself is not reproduced here for obvious
reasons.

\begin{table}[h]
  \centering
  \small
  \begin{tabular}{ll}
    \toprule
    \textbf{Parameter} & \textbf{Value} \\
    \midrule
    Channel index         & 0 (primary) \\
    PSK type              & Random, 256-bit \\
    PSK identical across nodes & Yes (verified via channel URL) \\
    Uplink/downlink enabled & No \\
    \bottomrule
  \end{tabular}
  \caption[Channel security configuration]{\label{tab:bijlage-security}Channel security configuration, identical on all three nodes.}
\end{table}